%! Dividing Polynomials; Remainder and Factor Theorems — Unit 4, Lesson 4.3 — Homework (Answer Key)
\documentclass[10pt]{article}
\usepackage{algebra2-article}
\usepackage{algebra2-key}   % pulls in -boxes; do NOT also load -boxes
\newcommand{\TallMath}[1]{$\displaystyle #1\rule[-1.4em]{0pt}{3.2em}$}

\begin{document}

\pageheader{Unit 4, Lesson 4.3}{Homework --- Answer Key}
\namedateperiod

\vspace{0.10in}

\begin{notesbox}{Practice}
Standard form and placeholders first, every time. When a remainder is not $0$, write the answer as
quotient $+\;\dfrac{\text{remainder}}{\text{divisor}}$.

\begin{enumerate}[leftmargin=1.9em, itemsep=11pt]
  \item \textbf{Monomial divisors.}
    \begin{enumerate}[label=(\alph*), leftmargin=1.8em, itemsep=5pt]
      \item $\dfrac{16x^5 - 24x^3 + 8x^2}{8x^2} = $ \ans{$2x^3 - 3x + 1$}
      \item $\dfrac{18x^3y^2 - 27x^2y^3}{9x^2y^2} = $ \ans{$2x - 3y$}
      \item $\dfrac{20x^4 + 5x^3 - 10x^2}{5x^2} = $ \ans{$4x^2 + x - 2$}
    \end{enumerate}

  \item \textbf{Long division.} Show the tableau on your own paper; record the results here.
    \begin{enumerate}[label=(\alph*), leftmargin=1.8em, itemsep=5pt]
      \item $(x^3 + 2x^2 - x - 2) \div (x - 1)$: \; quotient \ans{$x^2+3x+2$}, remainder \ans{$0$}
      \item $(2x^3 - 3x^2 + 4) \div (x - 2)$: \; quotient \ans{$2x^2+x+2$}, remainder \ans{$8$}

        \vspace{2pt}
        What placeholder did (b) need? \ans{$0x$}
      \item $(x^4 - 2x^3 - 7x^2 + 8x + 12) \div (x^2 - x - 6)$: \; quotient \ans{$x^2-x-2$},
        remainder \ans{$0$}

        \vspace{2pt}
        Why could you \emph{not} use synthetic division on (c)? \ans{The divisor is quadratic, not of the form $x-r$.}
    \end{enumerate}

  \item \textbf{Synthetic division.} Give the quotient and the remainder.
    \begin{enumerate}[label=(\alph*), leftmargin=1.8em, itemsep=5pt]
      \item $(x^3 - 7x + 6) \div (x - 1)$: \; \ans{$x^2+x-6$, remainder $0$}
      \item $(3x^3 + 8x^2 - x - 10) \div (x + 2)$: \; \ans{$3x^2+2x-5$, remainder $0$}
      \item $(x^4 - 16) \div (x - 2)$: \; \ans{$x^3+2x^2+4x+8$, remainder $0$}

        \vspace{2pt}
        In Lesson 4.2 you factored $x^4-16$ completely. Does your quotient here agree with that
        answer? Explain in one line. \ans{Yes --- $x^4-16=(x-2)(x+2)(x^2+4)$, and $(x+2)(x^2+4)=x^3+2x^2+4x+8$.}
    \end{enumerate}

  \item \textbf{Remainder Theorem --- no dividing.} Let $f(x) = x^3 - 2x^2 + 3x - 4$.
    \begin{enumerate}[label=(\alph*), leftmargin=1.8em, itemsep=5pt]
      \item Remainder on division by $(x-1)$: \ans{$-2$} \qquad by $(x+1)$: \ans{$-10$}
            \qquad by $(x-2)$: \ans{$2$}
      \item Are any of those three binomials a factor of $f$? \ans{No} \; How can you tell?
            \ans{None of the three remainders is $0$ (Factor Theorem).}
    \end{enumerate}

  \item \textbf{Factor Theorem --- test, divide, finish.} Let $p(x) = x^3 + 2x^2 - 11x - 12$.
    \begin{enumerate}[label=(\alph*), leftmargin=1.8em, itemsep=5pt]
      \item Show that $x = -1$ is a zero: $p(-1) = $ \ans{$0$}
      \item So \ans{$(x+1)$} is a factor. Divide it out: depressed polynomial \ans{$x^2+x-12$}
      \item Factor completely: $p(x) = $ \ans{$(x+1)(x+4)(x-3)$}
      \item List the zeros: \ans{$-1,\ -4,\ 3$}
    \end{enumerate}

  \item \textbf{Back to a graph you already read.} On Lesson 4.0's exit ticket you read the graph of
        $h(x) = x^3 - 3x - 2$ --- it touches the $x$-axis at $x = -1$ and crosses at $x = 2$.
    \begin{enumerate}[label=(\alph*), leftmargin=1.8em, itemsep=5pt]
      \item Confirm with the Factor Theorem: $h(-1) = $ \ans{$0$}
      \item Divide $h(x)$ by $(x+1)$: quotient \ans{$x^2-x-2$}
      \item Factor the quotient and write $h$ completely: \ans{$(x+1)^2(x-2)$}
      \item The zero $x=-1$ has multiplicity \ans{$2$}, which is why the graph
            \ans{touches} there instead of crossing.
    \end{enumerate}

  \item \textbf{Explain.} A classmate divides $f(x)$ by $(x - 5)$ and gets a remainder of $0$. In two
        or three sentences, tell them three things they now know about $f$ --- and say which
        statement in the division statement $P = D\cdot Q + R$ makes it true. \ansline{They know that $(x-5)$ is a factor of $f$, that $f(5)=0$, and that $(5,0)$ is an $x$-intercept of the graph of $f$ --- and that the quotient they just found is the other factor. It follows from $P = D\cdot Q + R$: with $R=0$ the statement collapses to $f(x)=(x-5)\,Q(x)$, a product, so substituting $x=5$ makes the first factor zero and therefore $f(5)=0$.}
    \end{enumerate}
\end{notesbox}

\begin{extensionbox}
\textbf{Run the theorems backward.}
\begin{enumerate}[label=(\alph*), leftmargin=1.8em, itemsep=5pt]
  \item Find the value of $k$ that makes $(x - 2)$ a factor of $f(x) = x^3 + kx^2 - 4x + 4$, then
        factor the resulting polynomial completely. \ans{$f(2)=8+4k-8+4=4k+4=0$, so $k=-1$ and $f(x)=x^3-x^2-4x+4=(x-2)(x+2)(x-1)$.}
  \item A rectangular prism has volume $V(x) = x^3 + 6x^2 + 11x + 6$ cubic units and height
        $(x + 1)$ units. Find an expression for the area of its base, then give the two base
        dimensions as binomials. \ans{Base area $=V\div(x+1)=x^2+5x+6=(x+2)(x+3)$, so the base is $(x+2)$ by $(x+3)$.}
  \item Check your answer to (b) at $x = 2$: compute $V(2)$ directly, and multiply your three
        dimensions at $x=2$. Do they agree? \ans{$V(2)=8+24+22+6=60$; the box is $3\times 4\times 5=60$ --- yes.}
  \item Synthetic division evaluates as well as it divides. Use it to find $f(3)$ for
        $f(x) = 2x^4 - 5x^3 + x - 7$, then verify by substituting $3$ directly. \ansline{Coefficients $2,\ -5,\ 0,\ 1,\ -7$ with $r=3$ give the bottom row $2,\ 1,\ 3,\ 10,\ 23$, so $f(3)=23$. Directly: $162-135+3-7=23$. \checkmark}
\end{enumerate}
\end{extensionbox}

\begin{spiralbox}
\textbf{Coming next --- Lesson 4.4: The Rational Root Theorem.} Today you could crack any polynomial
\emph{once somebody told you a zero to test}. Tomorrow you generate that list yourself: the Rational
Root Theorem turns the constant term and the leading coefficient into a short list of candidate
zeros, synthetic division tests them, and the depressed polynomial finishes the job. Notice already
that a zero can be a fraction --- $2x^3-3x^2-11x+6$ has the zero $\tfrac12$ --- so the list will need
more than whole numbers.
\end{spiralbox}

\begin{teachernote}
\textbf{Grading focus.} Items 4, 6, and 7 carry the standard; 1--3 are procedure. In item~2(b) the
missing $x$ term is the point --- a student who divides $2x^3-3x^2+4$ without the placeholder gets a
quotient of the wrong degree, and the error is visible at a glance. Item~2(c) is the only long
division in the set that synthetic division cannot do, which is the argument for keeping both tools.

\vspace{3pt}
\textbf{Item 3(c) is worth two minutes of class debrief:} dividing $x^4-16$ by $(x-2)$ returns
$x^3+2x^2+4x+8$, which by grouping is $(x+2)(x^2+4)$ --- exactly the factorization students produced
in Lesson~4.2 by a completely different route. Same answer, two roads.

\vspace{3pt}
\textbf{Item 6 closes the longest loop in the unit.} Students read this graph on Lesson~4.0's exit
ticket, expanded $(x+1)^2(x-2)$ into $x^3-3x-2$ for Lesson~4.1 homework, could not factor it back in
Lesson~4.2, and now recover the factored form by division. Say that out loud when you review it.
\textbf{Item 7} is read for reasoning: full credit needs all three equivalent statements (factor,
$f(5)=0$, $x$-intercept at $5$), not just the first.
\end{teachernote}

\end{document}
